% Figure 5: K-Scaling Experiment — Hybrid vs. Disjoint LinUCB

\begin{figure*}[t]
\centering
\includegraphics[width=0.95\textwidth]{figures/figure5_k_scaling.png}
\caption{\textbf{Family parameter sharing reduces convergence variance in the production router.}
Greedy holdout reward (750 prompts) at 25-step intervals during training on 533 prompts via the full Corralling router (warmup + tabula rasa experts).
Columns show portfolios of $K \in \{5, 10\}$ models; bands are 95\% CIs ($n{=}20$ seeds).
Family assignments use within-provider tetrachoric correlation ($r_{\text{tet}} \geq 0.6$) computed on the training set only.
Both portfolios include multi-member providers so shared families form at every $K$: 2 shared families at $K{=}5$ (4/5 models) and 3 at $K{=}10$ (6/10 models).
Hybrid produces consistently narrower CIs---$1.8\times$ at $K{=}5$ ($\pm 0.018$ vs.\ $\pm 0.033$) and $2.1\times$ at $K{=}10$ ($\pm 0.017$ vs.\ $\pm 0.035$)---though mean differences are not significant at the $0.05$ level ($p{=}0.078$ and $p{=}0.160$, respectively).
Arm-pull allocation is highly non-uniform (e.g.\ median 47, range 0--120 at $K{=}10$), confirming that per-model data scarcity is driven by bandit exploitation, not uniform assignment.}
\Description{Two-panel line plot comparing Hybrid and Disjoint holdout reward over training steps at K=5 and K=10, with 95 percent confidence bands showing Hybrid producing narrower intervals at both portfolio sizes.}
\label{fig:k_scaling}
\end{figure*}
